\section{Chapter 11}\label{chapter-11}

\subsection{Can We Trust an Autonomous
Operator?}\label{can-we-trust-an-autonomous-operator}

An authority chain inventory can show what has been assembled around a
system. It can show which model, connector, credential, service and
policy shape an action. It cannot answer the harder question:
\textbf{should this arrangement be allowed to act at all?}

That question is where trust begins.

The wrong way to ask it is personal. \emph{Can we trust the agent?}
sounds natural because people know how to trust a colleague. A colleague
can explain a choice, revise a view, accept a duty and answer for a
mistake. An agentic system does none of those things. It produces
outputs and effects through a model, context, tools, policies, services
and data supplied by an organisation.

The useful question is therefore narrower and more demanding:

\begin{quote}
\textbf{Can this organisation defend allowing this system to exercise
this class of autonomy in this context?}
\end{quote}

That is not a question of personality. It is a question of design.

\subsection{Trust Is Not a Feeling About a
System}\label{trust-is-not-a-feeling-about-a-system}

In January 2026, Singapore's Ministry of Digital Development and
Information published a Model AI Governance Framework for Agentic AI.
The framework did not propose that organisations should decide whether
to ``trust AI'' in the abstract. It described conditions for responsible
deployment: bound risks, meaningful human accountability, controls
across the lifecycle and informed end users.{[}1{]}

The distinction matters.

A system may perform reliably in one setting and still be indefensible
in another. A customer-service workflow might accurately classify a
request yet have no authority to alter an account. A hospital operations
workflow might efficiently prepare a staffing proposal yet have no
authority to override clinical escalation rules. An infrastructure
workflow might accurately identify a configuration change yet have no
authority to apply it to a live environment.

Performance is evidence about a system's behaviour under certain
conditions. Trust is a judgment about whether the whole arrangement may
be allowed to produce an effect.

Those are not the same question.

A high-performing model can sit inside a weak arrangement. It may
receive a permanent credential, accept untrusted instructions, invoke an
unreviewed connector and leave no usable record of what changed. A
cautious model can sit inside a strong arrangement. It may act only
inside a narrow mandate, through a temporary authority, with observable
evidence and a recovery path.

The first arrangement may look impressive. The second is more
trustworthy.

\subsubsection{Evidence Note: A Framework Is Not a
Certificate}\label{evidence-note-a-framework-is-not-a-certificate}

A published framework can identify useful conditions. It does not
certify a particular system as safe, reliable or trustworthy. The
relevant question remains local: \textbf{can an organisation demonstrate
that its own autonomy arrangement meets the conditions it claims to
meet?}

\subsection{The Defensible Trust
Proposition}\label{the-defensible-trust-proposition}

An organisation needs a sentence it can test.

For example:

\begin{quote}
\textbf{The invoice-reconciliation workflow may match approved purchase
orders to received invoices and prepare payment exceptions for review.
It may not release funds, create vendors, change payment instructions or
use credentials outside the finance environment. Its actions expire at
the end of the reconciliation window and are reconstructable from the
case record.}
\end{quote}

This is not a slogan. It is a \textbf{defensible trust proposition}.

A defensible trust proposition identifies an action class, a context, an
authority boundary, a source of evidence and a recovery condition. It
can be challenged before an incident, not merely narrated after one.

The proposition has five parts.

\subsubsection{1. Mandate}\label{mandate}

The mandate names what the system is allowed to pursue. It should be
expressed as an outcome or a bounded action class, not as a vague
aspiration such as ``help the team'' or ``improve operations.''

``Reconcile approved invoices'' is a mandate. ``Assist finance'' is not.

``Prepare a deployment plan from approved changes'' is a mandate.
``Operate production'' is not.

A mandate provides the first answer to the question of purpose. Without
it, a system's apparent usefulness can expand into a collection of
unexamined tasks. An organisation then discovers its actual policy only
after an effect has occurred.

\subsubsection{2. Authority}\label{authority}

Authority names the route through which the system can act. It includes
the identity it uses, the tools and data it can reach, the environment
in which those tools operate and the scope that limits each call.

A mandate without authority is only a request. Authority without a
mandate is an open-ended risk.

This distinction remains durable even as model capabilities change. A
model may become more able to plan, retrieve, write code or coordinate
tools. None of those improvements tell an organisation what the system
is entitled to alter. The entitlement must be designed elsewhere.

The NIST AI Agent Standards Initiative, announced in 2026, reflects the
same shift from a model-only view toward secure, interoperable systems
that can take autonomous actions.{[}2{]} The important implication is
not that a standard can solve trust. It is that trust has to be
expressed across the components that give an action its reach.

\subsubsection{3. Evidence}\label{evidence}

Evidence is what allows the organisation to distinguish a claim from a
reconstruction.

A useful record does more than state that an action happened. It should
connect a mandate, a policy version, the relevant state, the identity
used, the tool invoked, the action produced, the result observed and the
recovery decision where one was needed.

This does not require an organisation to record every internal model
token or to pretend that every output has a simple explanation. It
requires enough evidence to answer a practical question: \textbf{what
would a reasonable reviewer need to see in order to decide whether the
system acted within its arrangement?}

That question applies outside technology teams. A claims workflow needs
to show which policy and documents supported a recommendation. A
customer-service workflow needs to show which account scope and
communication channel it used. A laboratory operations workflow needs to
show which dataset, procedure and escalation rule shaped an action. The
mechanism changes. The burden of explanation does not.

\subsubsection{4. Recovery}\label{recovery}

No arrangement is trustworthy because it claims it will never fail.

A trustworthy arrangement expects that a tool can change, a data source
can be wrong, a connector can behave unexpectedly or an action can reach
farther than intended. Recovery therefore belongs in the trust
proposition from the beginning.

Recovery has at least four questions. What stops the action? What can be
reversed? What must be contained? Who decides that the arrangement may
resume?

The answer may be technical, such as a scoped rollback or a credential
revocation. It may be operational, such as pausing a workflow and moving
work to a review queue. It may be organisational, such as escalating a
material exception to a named owner. The form is less important than the
fact that the path exists before the pressure of an incident.

\subsubsection{5. Accountability}\label{accountability}

Accountability is not the final name written at the bottom of a diagram.
It is the ability to change the arrangement and answer for its effects.

A person who can receive an alert but cannot narrow scope, suspend a
connector, change a policy or stop a workflow is not meaningfully
accountable. They are a spectator with a notification.

Meaningful accountability requires a decision right. The accountable
owner must be able to alter the mandate, authority, evidence requirement
or recovery condition when the arrangement no longer deserves its
autonomy.

This is why a system must not be described as a digital employee. It
cannot hold a decision right, accept public responsibility or revise the
organisation's purpose. It can execute within an arrangement that people
and institutions remain responsible for designing.

\subsection{Control Question}\label{control-question}

\begin{quote}
\textbf{If this system produced an unacceptable effect tomorrow, could
the organisation show its mandate, authority, evidence, recovery path
and accountable owner without constructing the answer after the fact?}

If the answer is no, the organisation has not yet built a trust
arrangement. It has built a hope arrangement.
\end{quote}

\subsection{Trust Has a Scope and an
Expiry}\label{trust-has-a-scope-and-an-expiry}

Trust becomes misleading when it is spoken of as a permanent property.

An organisation may be able to defend a system's autonomy for one action
class and not another. It may defend it in a test environment and not in
production. It may defend it while a policy, tool version and data
source remain unchanged, but not after a connector is replaced or a
credential is broadened.

A trust proposition should therefore have a scope and an expiry
condition.

The scope identifies where the proposition holds. The expiry condition
identifies what change requires the proposition to be reviewed. A new
model release, a new tool, an expanded data source, a different
environment or a changed owner can all invalidate an earlier assurance
claim.

This does not make autonomy impractical. It makes it honest.

The safest systems are not the systems that never change. They are the
systems whose changes force an organisation to decide whether the
autonomy arrangement still deserves to exist.

\subsection{From Control Lists to Trust
Architecture}\label{from-control-lists-to-trust-architecture}

Many organisations begin with a list: use role-based access, keep logs,
require review, test the model, train users. Each item can be sensible.
None is enough alone.

A list treats trust as the sum of controls. A trust architecture treats
it as a relationship among conditions.

An access control is meaningful only when it enforces a mandate. A log
is meaningful only when it can establish whether an action stayed inside
its authority. A review is meaningful only when it can change an
outcome. A test is meaningful only when it informs a decision about
scope or deployment. Training is meaningful only when users know when
and how to challenge the system's apparent confidence.

The architecture is what makes the components mutually intelligible.

Singapore's 2026 framework is useful because it pairs bounds on agent
powers with meaningful human accountability, lifecycle controls and user
responsibility.{[}1{]} It does not reduce trust to a single dashboard or
a single approval point. It treats deployment as an organisational
arrangement that must remain governable over time.

The details of the framework may evolve. The underlying idea will not:
autonomy becomes defensible when an organisation can explain why a
system may act, what constrains it, how it is observed, how it is
stopped and who can change those conditions.

\subsection{Assurance Is a Practice, Not a
Verdict}\label{assurance-is-a-practice-not-a-verdict}

A trust proposition is not completed when a system is deployed. It has
to remain open to challenge.

This is where assurance differs from a launch review. A launch review
asks whether the arrangement appears ready at a moment in time.
Assurance asks whether the organisation can continue to justify its
claim as the environment changes. The distinction becomes important the
first time a policy is edited, a connector is upgraded, an owner changes
teams or an apparently minor workflow is reused for a more consequential
task.

A system may still work after such a change. That is not enough. The
organisation must decide whether the proposition under which it was
allowed to act still holds.

NIST's risk-management work is useful here not because it supplies a
permanent score for trust, but because it treats management as a
continuing practice of governing, mapping, measuring and managing
risk.{[}3{]} An autonomous workflow needs an equivalent rhythm. It
should be reviewed when its authority changes, when evidence no longer
supports the claim, or when a failure reveals that recovery was less
available than the design assumed.

This makes trust inconvenient in a productive way. It prevents an
organisation from turning a successful pilot into a silent entitlement.

\subsubsection{A Trust Proposition Can Fail Without a
Failure}\label{a-trust-proposition-can-fail-without-a-failure}

A system does not need to produce harm before its trust proposition
becomes invalid.

Suppose a workflow has been allowed to reconcile invoices inside one
business unit. Its mandate is clear, its authority is bounded and its
outputs are reviewed. A later change gives the workflow access to a
shared vendor directory. No invoice has been paid incorrectly. No alert
has fired. Yet the authority condition has changed. The original claim
was about a workflow operating within a limited environment. The revised
arrangement may require a new claim, a new review and possibly a
narrower scope.

The same pattern appears in clinical administration, public services and
customer operations. A scheduling system can begin by proposing
appointments, then gain access to a patient communication channel, then
be connected to an exception queue. Each connection may seem useful on
its own. Together they change what the arrangement can cause. The
relevant question is not whether the system has become more intelligent.
It is whether the organisation has expanded autonomy without renewing
its justification.

This is the practical meaning of trust with an expiry condition. The
arrangement does not lose trust because it changed. It loses the right
to rely on an old trust claim after it changed.

\subsubsection{The Difference Between Monitoring and Meaningful
Review}\label{the-difference-between-monitoring-and-meaningful-review}

Monitoring reports what a system is doing. Meaningful review tests
whether the organisation should still permit it.

A dashboard can show a high completion rate while hiding an authority
problem. A low error rate can show that a workflow is operating
consistently, but not whether it is operating within a legitimate
purpose. An alert can identify an unusual tool call, but not determine
whether the mandate itself was too broad.

Meaningful review has to include someone with the power to question the
proposition, not merely the power to observe it. The reviewer needs to
ask whether the action class is still appropriate, whether the evidence
is still sufficient, whether the recovery path still works and whether
the accountable owner can still act.

That role need not be a committee in every case. A low-impact internal
workflow may require a lightweight operational review. A system that can
alter customer records, make payments, affect a safety-sensitive process
or change production infrastructure may require a more formal challenge.
The principle is stable: \textbf{the strength of review should follow
the consequence of autonomy, not the novelty of the technology.}

\subsection{The Human Remainder}\label{the-human-remainder}

The human remainder is not the work left over after automation.

It is the authority an organisation must keep because it cannot be
transferred without changing the organisation's own responsibility.
People define the purpose of a workflow. They decide which harms are
acceptable or unacceptable. They choose when a recovery threshold has
been crossed. They decide whether evidence is sufficient to resume a
suspended system.

These are not gaps in current technology. They are organisational
duties.

A system may calculate a route, rank a queue, prepare a response, match
a document or execute a bounded change. It cannot determine what an
organisation is entitled to optimise at the expense of whom. That is a
question of purpose and legitimacy, not of inference.
